\documentclass[a4paper]{article}
\usepackage{amsfonts}
\usepackage{amsmath}
\usepackage{amssymb}
\usepackage{graphicx}     

\begin{document}
  
\noindent \large Auteur: Michael Livchits \\
\noindent \large Studierichting: Informatica\\
\noindent \large Oefeningenreeks 3D nr. 4\\

\medskip

\normalsize

Door de eerste stelling weten wij dat:\\

$A(x) = (x-1)(x+1)Q(x)$\\

Door stelling 3 weten wij dat $A'(1) = 0$:\\

Dit betekent dus dat zowel $A(1) = 0$ en $A'(1) = 0$\\

$A(x) = (x-1)^2(x+1)(x+b)$ met $b \in \mathbb{R}$\\

Nu kunnen wij met stelling 2, $b$ vinden:\\

$A(x)=(x-1)^2(x+1)(x+b)$\\

$A\left(-\dfrac{1}{2}\right)=\left(-\dfrac{1}{2}-1\right)^2\left(-\dfrac{1}{2}+1\right)\left(-\dfrac{1}{2}+b\right)$\\

$-\dfrac{27}{16} = \left(-\dfrac{3}{2}\right)^2\left(\dfrac{1}{2}\right)\left(b-\dfrac{1}{2}\right)$\\

$-\dfrac{27}{16} = \dfrac{9}{8}\left(b-\dfrac{1}{2}\right)$\\

$-\dfrac{3}{2} = b-\dfrac{1}{2}$\\

$b=-1$\\

Met andere woorden laatste nulpunt is $(x+b) \Rightarrow (x-1)$\\

$A(x) = (x-1)^3(x+1)$\\

$A(x) = x^4-2x^3+2x-1$\\

\begin{thebibliography}{999}
%\bibitem{a} Referentie
\end{thebibliography}
\end{document}